\chapter{Dental Exam}
\section*{What Is This Test?} A dental exam is a clinical assessment of the teeth, gums, mouth, jaw, and bite to detect decay, gum disease, infection, oral lesions, and other problems.
\section*{Alternative Names} Oral examination; dental checkup; comprehensive dental examination.
\section*{Principle} Inspection, palpation, periodontal measurements, and targeted radiographs assess visible and hidden structures.
\section*{Why This Test Is Ordered} Routine prevention, tooth pain, swelling, bleeding gums, trauma, lesions, or preparation for dental treatment.
\section*{Primary vs Secondary Test} History and examination are primary; X-rays, cultures, biopsy, and blood tests are selected secondary studies.
\section*{How This Test Is Performed} The dentist inspects the mouth, checks teeth and gums, measures pockets, evaluates the bite, and orders imaging when indicated.
\section*{What Preparation Is Needed} No special preparation is usually required. Bring medication, dental, and medical history.
\section*{Understanding Results} Findings may include caries, gingivitis, periodontitis, infection, abnormal growth, or a need for preventive care.
\section*{Follow-up Tests} Dental radiographs, periodontal charting, pulp testing, culture, biopsy, or referral may be needed.
\section*{References} \begin{enumerate}\item MedlinePlus. Dental Exam.\item American Dental Association. Oral health examination guidance.\end{enumerate}
\clearpage\section*{Understanding Results and Follow-up}
The laboratory report should identify the specimen, method, units, reference interval or decision limit, and whether the result is positive, negative, abnormal, or inconclusive. Reference intervals vary with age, sex, pregnancy, specimen type, assay, and laboratory; a result outside a range is not automatically a diagnosis, and a result within a range does not always exclude disease. Discuss unexpected results with the ordering clinician, who can decide whether repeat or confirmatory testing, treatment, or referral is appropriate. Do not change medicines or delay urgent care based on an isolated result.
\section*{Regulatory Status and Authorization Date}This chapter describes a test category, not a specific commercial product. FDA, CDSCO, EU IVDR, and MHRA approval or registration dates therefore cannot be assigned without the assay manufacturer, product name, instrument, and intended use. For laboratory-developed tests, an individual agency approval date may not exist. See the Preface for the interpretation of regulatory status.
\clearpage
